% =====================================================================
%  ch03_rtl.tex -- 第 3 章 中值滤波与电压换算模块 rtl.v
% =====================================================================
\section{中值滤波与电压换算模块 \texttt{rtl}（\file{src/rtl.v}）}

\subsection{文件信息}

\begin{center}
\small
\begin{tabular}{ll}
\toprule
文件路径 & \file{FPGA/src/rtl.v}（注释后 143 行）\\
模块 / 实例 & \texttt{rtl} / \sig{adc_inst}（在 \texttt{adda} 中）\\
功能 & (1) 用 1024 点平均值标定 ADC 零点；(2) 换算正/负向电压\\
时钟 & \sig{clk}=50MHz，\sig{rst_n} 低有效\\
实现代价 & 2 个 30bit 除法器、1 个 14bit 除法、若干乘加\\
\bottomrule
\end{tabular}
\end{center}

\subsection{第 1 层：模块接口}

\begin{table}[htbp]
\centering
\small
\caption{rtl 端口}
\begin{tabular}{llll}
\toprule
端口 & 方向 & 位宽 & 说明\\
\midrule
\sig{clk} & input & 1 & 系统时钟 50MHz\\
\sig{rst_n} & input & 1 & 异步复位，低有效\\
\sig{ad_data} & input & 10 & ADC 码值（0$\sim$1023 无符号）\\
\sig{sign} & output & 1 & 1 表示码值低于零点（负电压）\\
\sig{volt} & output & 10 & 电压值，单位 0.01V（$100\times$ 标度）\\
\bottomrule
\end{tabular}
\end{table}

\subsection{第 2 层：内部信号}

\begin{table}[htbp]
\centering
\small
\begin{tabular}{llp{8cm}}
\toprule
信号 & 位宽 & 作用\\
\midrule
\sig{clk_sample} & 1 & 预留，未使用\\
\sig{median_en} & 1 & 标定完成标志，置 1 后保持\\
\sig{cnt_median} & 11 & 累加计数 0$\sim$1024\\
\sig{data_sum_m} & 30 & 1024 点累加和\\
\sig{data_median} & 14 & 零点（平均值）\\
\sig{volt_reg} & 31 & 电压中间量\\
\sig{data_p} & 31 & 正向每 LSB 电压刻度（定点 $2^{13}\times100$ 倍）\\
\sig{data_n} & 31 & 负向每 LSB 电压刻度\\
\bottomrule
\end{tabular}
\end{table}

\subsection{第 3 层：语句块解析}

\subsubsection{语句 1：符号判断（46--51 行）}

\begin{lstlisting}[style=vsnip]
assign sign=(ad_data<data_median)? 1'b1:1'b0;
//assign sign=1'b1;
\end{lstlisting}

当前码值低于标定零点即判为负，输出 \sig{sign} 给 \sig{seg} 决定是否显示负号。
注释掉的 \sig{assign sign=1'b1} 是调试时强制显示正号的开关。

\subsubsection{块 1：中值使能（53--65 行）}

\begin{lstlisting}[style=vsnip]
always @(posedge clk or negedge rst_n)
    if(!rst_n)
        median_en<=1'b0;
    else if(cnt_median==CNT_DATA_MAX)
        median_en<=1'b1;
    else
        median_en<=median_en;
\end{lstlisting}

\sig{cnt_median} 数到 1024 时把 \sig{median_en} 置 1 并保持。由于块 2 在
\sig{median_en}=1 后停止计数，\sig{cnt_median} 永远停在 1024，所以
\sig{median_en} 一旦置位就不会回零——这是一次性的“上电自校准”。

\subsubsection{块 2：累加计数器（67--76 行）}

\begin{lstlisting}[style=vsnip]
always @(posedge clk or negedge rst_n)
    if(!rst_n)
        cnt_median<=11'b0;
    else if(!median_en)
        cnt_median<=cnt_median+1'b1;
\end{lstlisting}

注意此处\textbf{没有 else 分支}：当 \sig{median_en}=1 时 \sig{cnt_median} 保持不变
（推断为寄存器，不会产生锁存器，因为复位分支存在）。这个“计数到上限后自锁”
的写法与块 1 配合完成单次标定。

\subsubsection{块 3：1024 点累加（78--90 行）}

\begin{lstlisting}[style=vsnip]
always @(posedge clk or negedge rst_n)
    if(!rst_n)
        data_sum_m<=30'b0;
    else if(cnt_median==CNT_DATA_MAX)
        data_sum_m<=30'b0;
    else
        data_sum_m<=data_sum_m+ad_data;
\end{lstlisting}

在计数未满 1024 时把 \sig{ad_data} 累加；计满的那一拍清零。最大累加和
$1023\times1024=1\,047\,552<2^{20}$，用 30bit 有充足余量。

\subsubsection{块 4：求平均（92--103 行）}

\begin{lstlisting}[style=vsnip]
always @(posedge clk or negedge rst_n)
    if(!rst_n)
        data_median<=12'b0;
    else if(cnt_median==CNT_DATA_MAX)
        data_median<=data_sum_m/CNT_DATA_MAX;
    else
        data_median<=data_median;
\end{lstlisting}

累加和除以 1024。综合工具会把“除以 2 的整数次幂”优化成右移 10 位，不消耗除法器。
由于非阻塞赋值，\sig{data_median} 取到的是清零\textbf{之前}的累加和，结果正确。
得到的平均值即 ADC 的零点（板级直流偏置），典型值约 512。

\subsubsection{语句 2：正/负向电压刻度（105--115 行）}

\begin{lstlisting}[style=vsnip]
assign data_p=(median_en==1'b1)?3276800/((1024-data_median)*2):0;
assign data_n=(median_en==1'b1)?3276800/((data_median+1)*2):0;
\end{lstlisting}

把满量程 $\pm2.0$V 按“零点两侧各自的码值区间”等分，得到每 LSB 的电压（定点数）。
$3276800=2^{13}\times400$，于是
\begin{align}
\text{data\_p} &= \frac{2^{13}\times400}{2(1024-m)} = \frac{2^{13}\times200}{1024-m},\\
\text{data\_n} &= \frac{2^{13}\times400}{2(m+1)} = \frac{2^{13}\times200}{m+1},
\end{align}
其中 $m=\text{data\_median}$。随后块 5 再右移 13 位，恰好抵消 $2^{13}$，得到
$100\times V$ 的定点电压。两个除法器都是 30bit 组合除法，需要多级移位减法，
是本模块 LUT 的主要来源。

\subsubsection{块 5：电压换算（117--136 行）}

\begin{lstlisting}[style=vsnip]
always @(posedge clk or negedge rst_n)
    if(!rst_n)
        volt_reg<=1'b0;
    else if(median_en)
        if((ad_data>(data_median-3)) && (ad_data<(data_median+3)))
            volt_reg<=1'b0;
        else if(ad_data<data_median)
            volt_reg<=(data_n*(data_median-ad_data)) >> 13;
        else if(ad_data>data_median)
            volt_reg<=(data_p*(ad_data-data_median)) >> 13;
    else
        volt_reg<=1'b0;
\end{lstlisting}

三条规则：
\begin{enumerate}[leftmargin=2em,itemsep=2pt]
  \item \textbf{死区}：$|ad\_data-m|\le 2$ 时输出 0，避免零点附近显示末位跳动；
  \item \textbf{负向}：\sig{data_n}$\times(m-ad\_data)$ 后右移 13；
  \item \textbf{正向}：\sig{data_p}$\times(ad\_data-m)$ 后右移 13。
\end{enumerate}

\sig{volt} 的有效范围是 0$\sim$1023（对应 0$\sim$10.23V 显示量程），而实际物理量程是
$\pm2.0$V，因此数值上 $\text{volt}=100\times|\text{电压}|$，例如 AD 满码时约 199（1.99V）。

\begin{tipbox}
\textbf{数值验算}：$m=512$，$ad\_data=1023$ 时
$\text{data\_p}=2^{13}\times200/512=3200$，
$\text{volt}=(3200\times511)>>13=199$，即 1.99V；\sig{ad_data}=0 时同理得 199。
与 $\pm2$V 满量程、10bit 分辨率（约 3.9mV/LSB）一致。
\end{tipbox}

\subsubsection{语句 3：输出（138--141 行）}

\begin{lstlisting}[style=vsnip]
assign volt=volt_reg;
\end{lstlisting}

31bit 中间量直接连到 10bit 端口，综合时会隐式截断（取低 10 位）。由于设计保证
结果不超过 1023，实际不会截断出错。

\subsection{本章小结与提示}

\begin{itemize}[leftmargin=2em,itemsep=2pt]
\item 本模块的核心思想是“\textbf{上电自校准 + 分段线性换算}”，用一个平均值把 ADC 的
      直流量化偏置消掉，再把码值差换算成电压；
\item 中值只在复位后 20.5$\mu$s 内计算一次，之后全程锁存；若温度漂移大，可改成
      “周期性重标定”（例如每 1s 重算一次）；
\item 两个除法器占用了较多 LUT，若资源紧张，可把 $1/(1024-m)$ 改为
      “查表 + 线性插值”或直接用移位近似；
\item 输出 \sig{volt} 在顶层没有负载，属“演示性输出”，实际综合时会被大幅裁剪。
\end{itemize}

\subsection{完整注释源码}

\lstinputlisting[style=vcode]{../FPGA/src/rtl.v}
